% [VERIFY] abstract claims trace to: Book.md abstract+claims_summary; C02,C03; Table 2,3,5,8
Deeper neural networks are more difficult to train. As feed-forward convolutional
networks for image recognition grow deeper, a counter-intuitive \degprob emerges:
adding layers to a suitably deep \plainnet increases not only test error but also
\emph{training} error, even with batch normalization and competent initialization,
so the deeper model is no easier to optimize than its shallower counterpart. Here,
we present a residual learning framework that eases the training of substantially
deeper networks by reformulating each block to fit a residual mapping
$\Fmap(\xin) := \Hmap(\xin) - \xin$ and recovering the desired mapping through a
parameter-free identity shortcut $\Fmap(\xin) + \xin$. We provide comprehensive
empirical evidence that these residual networks are easier to optimize and gain
accuracy from considerably increased depth: on ImageNet, residual nets reach a depth
of up to 152 layers---eight times deeper than VGG nets yet of lower complexity---and a
six-model ensemble attains 3.57\% top-5 error on the test set, winning the first place
in the ILSVRC 2015 classification task. We further analyze residual nets on CIFAR-10
with depths up to 110 and a 1202-layer stress test, and show that the learned
representations transfer to object detection, where replacing a VGG-16 backbone with a
101-layer residual net yields a 28\% relative improvement on COCO. The depth of
representations is of central importance for visual recognition.
